\documentclass[11pt]{article}

\usepackage[a4paper,margin=1in]{geometry}
\usepackage{amsmath,amssymb,amsthm,mathtools}
\usepackage[T1]{fontenc}
\usepackage{lmodern}
\usepackage{microtype}

\newtheorem{theorem}{Theorem}[section]
\newtheorem{lemma}[theorem]{Lemma}
\newtheorem{proposition}[theorem]{Proposition}
\newtheorem{corollary}[theorem]{Corollary}
\newtheorem{remark}[theorem]{Remark}
\newtheorem{conjecture}[theorem]{Conjecture}

\newcommand{\C}{\mathbb C}
\newcommand{\R}{\mathbb R}
\newcommand{\eps}{\varepsilon}

\title{Current Status of the Real-Slice Connectedness Problem\\
for the Clean OBC Hatano--Nelson Family}
\author{}
\date{}

\begin{document}
\maketitle

\section{Setup}

Fix \(0<a<b\), and let
\[
A_n=\operatorname{tridiag}(a,0,b)\in\R^{n\times n}.
\]
For \(z=x+iy\in\C\), define
\[
M_n(x,y)\coloneqq (zI-A_n)^*(zI-A_n),
\qquad
\lambda_*(x,y)\coloneqq \lambda_{\min}(M_n(x,y)).
\]
Then
\[
z\in\sigma_\eps(A_n)
\quad\Longleftrightarrow\quad
\lambda_*(x,y)<\eps^2.
\]
The connectedness-equivalence problem is to prove
\[
\sigma_\eps(A_n)\ \text{is connected}
\quad\Longleftrightarrow\quad
\sigma_\eps(A_n)\cap\R\ \text{is connected},
\]
or, equivalently, that for each fixed \(x\in\R\), the vertical slice
\[
\{y\in\R:\ \lambda_*(x,y)<\eps^2\}
\]
is either empty or an interval containing \(0\).

Write
\[
q\coloneqq b-a>0,
\qquad
\eta\coloneqq ab,
\qquad
c\coloneqq (a+b)x-iqy.
\]

\section{A topological reduction}

\begin{lemma}\label{lem:topological}
Let \(\Omega\subset\C\) be open. Assume that for every \(x\in\R\), the vertical section
\[
V_x(\Omega)\coloneqq \{y\in\R:\ x+iy\in\Omega\}
\]
is either empty or an interval containing \(0\). Then the connected components of
\(\Omega\) are in one-to-one correspondence with the connected components of
\(\Omega\cap\R\). In particular,
\[
\Omega\ \text{is connected}
\quad\Longleftrightarrow\quad
\Omega\cap\R\ \text{is connected}.
\]
\end{lemma}

\begin{proof}
This is the standard vertical-strip argument: if \(I\) is a connected component of
\(\Omega\cap\R\), then
\[
\Omega_I\coloneqq \{z\in\Omega:\ \Re z\in I\}
\]
is path connected because every nonempty vertical slice contains \(0\); and distinct
components of \(\Omega\cap\R\) are separated by a real interval whose entire vertical strip
is disjoint from \(\Omega\). Thus the \(\Omega_I\) are exactly the connected components of
\(\Omega\).
\end{proof}

\section{Exact odd/even block reduction}

Fix a threshold \(\Lambda\in\R\), and set
\[
\widehat H_n(\Lambda;x,y)\coloneqq M_n(x,y)-\Lambda I_n.
\]

\subsection*{Even case \(n=2m\)}

Define
\[
w_k\coloneqq \binom{u_{2k-1}}{u_{2k}},
\qquad
1\le k\le m.
\]
Then the equation
\[
\widehat H_{2m}(\Lambda;x,y)u=0
\]
is equivalent to the \(2\times2\) block Jacobi system
\[
B_Lw_1+Aw_2=0,
\]
\[
A^*w_{k-1}+Bw_k+Aw_{k+1}=0,
\qquad 2\le k\le m-1,
\]
\[
A^*w_{m-1}+B_Rw_m=0,
\]
where
\[
A=
\begin{pmatrix}
\eta & 0\\
-c & \eta
\end{pmatrix},
\]
\[
B_L=
\begin{pmatrix}
x^2+y^2+a^2-\Lambda & -c\\
-\bar c & x^2+y^2+a^2+b^2-\Lambda
\end{pmatrix},
\]
\[
B=
\begin{pmatrix}
x^2+y^2+a^2+b^2-\Lambda & -c\\
-\bar c & x^2+y^2+a^2+b^2-\Lambda
\end{pmatrix},
\]
\[
B_R=
\begin{pmatrix}
x^2+y^2+a^2+b^2-\Lambda & -c\\
-\bar c & x^2+y^2+b^2-\Lambda
\end{pmatrix}.
\]

\subsection*{Odd case \(n=2m+1\)}

Define
\[
w_k\coloneqq \binom{u_{2k-1}}{u_{2k}},
\qquad
1\le k\le m,
\qquad
t\coloneqq u_{2m+1}.
\]
Then the equation
\[
\widehat H_{2m+1}(\Lambda;x,y)u=0
\]
is equivalent to
\[
B_Lw_1+Aw_2=0,
\]
\[
A^*w_{k-1}+Bw_k+Aw_{k+1}=0,
\qquad 2\le k\le m-1,
\]
\[
A^*w_{m-1}+Bw_m+d_0\,t=0,
\]
\[
d_0^*w_m+e\,t=0,
\]
where
\[
d_0\coloneqq \binom{\eta}{-c},
\qquad
e\coloneqq x^2+y^2+b^2-\Lambda.
\]

\section{Tail Schur complements and the exact \(2\times2\) invariant}

The finite-dimensional spectral problem is encoded by the first tail Schur complement.

\begin{proposition}[Parity Weyl matrix / tail Schur complement]\label{prop:schur}
Let \(n\ge 2\).

\smallskip
\noindent
\textup{(i) Even case \(n=2m\).} Define recursively
\[
\Sigma_m\coloneqq B_R,
\qquad
\Sigma_k\coloneqq B_k-A\Sigma_{k+1}^{-1}A^*,
\qquad
k=m-1,\dots,1,
\]
where \(B_1=B_L\) and \(B_k=B\) for \(2\le k\le m-1\).

\smallskip
\noindent
\textup{(ii) Odd case \(n=2m+1\).} On the open set \(e\neq 0\), define
\[
\Sigma_m\coloneqq B-d_0e^{-1}d_0^*,
\qquad
\Sigma_k\coloneqq B_k-A\Sigma_{k+1}^{-1}A^*,
\qquad
k=m-1,\dots,1,
\]
again with \(B_1=B_L\) and \(B_k=B\) for \(2\le k\le m-1\).

\smallskip
Then in both parities,
\[
\det \widehat H_n(\Lambda;x,y)=0
\quad\Longleftrightarrow\quad
\det \Sigma_1(\Lambda;x,y)=0,
\]
up to the obvious harmless factor \(e\) in the odd case. In particular, the full
finite-\(n\) problem is encoded by the \(2\times2\) Hermitian matrix \(\Sigma_1\).
\end{proposition}

\begin{remark}
This \(2\times2\) Hermitian matrix \(\Sigma_1\) is the exact invariant object discovered so far:
it is the parity Weyl matrix (equivalently, the finite parity Dirichlet-to-Neumann map).
\end{remark}

\section{Scalar Riccati form and a first sign invariant}

Write
\[
\Sigma_k=
\begin{pmatrix}
u_k & \beta_k\\
\bar\beta_k & v_k
\end{pmatrix},
\qquad
\Delta_k\coloneqq u_kv_k-|\beta_k|^2.
\]
Then the interior Riccati step is
\[
u_k=\alpha_k-\frac{\eta^2v_{k+1}}{\Delta_{k+1}},
\]
\[
\beta_k=-c+\frac{\eta(\eta\beta_{k+1}+v_{k+1}\bar c)}{\Delta_{k+1}},
\]
\[
v_k=\delta_k-\frac{
\eta^2u_{k+1}
+\eta(\beta_{k+1}c+\bar\beta_{k+1}\bar c)
+|c|^2v_{k+1}
}{\Delta_{k+1}},
\]
with the obvious endpoint values
\[
(\alpha_1,\delta_1)=
(x^2+y^2+a^2-\Lambda,\ x^2+y^2+a^2+b^2-\Lambda),
\]
\[
(\alpha_k,\delta_k)=
(x^2+y^2+a^2+b^2-\Lambda,\ x^2+y^2+a^2+b^2-\Lambda)
\quad (2\le k\le m-1),
\]
and terminal seed \(\Sigma_m\) as above.

Set
\[
\tau_k\coloneqq \Im \beta_k.
\]
Then the imaginary part recursion is exact:
\[
\tau_k
=
qy+\frac{\eta(\eta\tau_{k+1}+qy\,v_{k+1})}{\Delta_{k+1}}.
\]
Therefore
\[
\tau_k>qy
\qquad\text{for every interior }k.
\]

\section{Projective upper half-plane invariance}

At a boundary point
\[
\Lambda=\lambda_*(x,y),
\qquad y>0,
\]
the tail matrices satisfy
\[
\Sigma_k\succ 0 \quad (k\ge 2),
\qquad
\Sigma_1\ge 0,
\qquad
\det\Sigma_1=0.
\]
Choose
\[
w_1\in \ker \Sigma_1\setminus\{0\},
\qquad
w_{k+1}\coloneqq -\Sigma_{k+1}^{-1}A^*w_k
\quad (k\ge 1),
\]
so that \(w_k\) is the block nullchain of the boundary singular vector.

For
\[
r_k\coloneqq \frac{w_{k,2}}{w_{k,1}},
\]
the exact projective recursion is
\[
r_{k+1}
=
\Phi_{k+1}(r_k)
\coloneqq
\frac{\eta u_{k+1}r_k-\bar\beta_{k+1}(\eta-\bar c\,r_k)}
{v_{k+1}(\eta-\bar c\,r_k)-\eta\beta_{k+1}r_k}.
\]

\begin{proposition}[Upper half-plane invariant]\label{prop:H}
At every boundary point with \(y>0\), one has
\[
r_k\in \mathbb H\coloneqq \{z\in\C:\ \Im z>0\}
\qquad (1\le k\le m).
\]
Equivalently,
\[
j_k\coloneqq \Im(\bar w_{k,1}w_{k,2})>0
\qquad (1\le k\le m).
\]
\end{proposition}

\begin{proof}[Proof sketch]
Because \(\Sigma_1\ge 0\), \(\det\Sigma_1=0\), and \(\tau_1>0\), one may choose a kernel vector
with \(r_1\in\mathbb H\). A direct calculation shows that if
\[
\Sigma=
\begin{pmatrix}
u & \beta\\
\bar\beta & v
\end{pmatrix},
\qquad
\beta=b+i\tau,
\qquad
\Delta=uv-b^2-\tau^2>0,
\qquad
\tau>qy,
\]
then \(\Phi(\mathbb H)\subset\mathbb H\). Applying this inductively yields
\(r_k\in\mathbb H\) for all \(k\).
\end{proof}

\section{A complete positive cocycle theorem for the shifted branch}

The next statement is fully proved.

\begin{proposition}[Positive cocycle for the shifted branch]\label{prop:Ppositive}
At a boundary point \(\Lambda=\lambda_*(x,y)\) with \(y>0\), define
\[
S_k\coloneqq -\partial_\Lambda \Sigma_k,
\qquad
P_k\coloneqq 2y\,S_k-\partial_y\Sigma_k.
\]
Then
\[
S_k\succ 0
\qquad\text{for all }k,
\]
and along the nullchain \(w_{k+1}=-\Sigma_{k+1}^{-1}A^*w_k\) one has
\[
w_1^*P_1w_1>0.
\]
More precisely, for each interior step \(k=1,\dots,m-1\),
\[
P_k=\Pi_k+R_kP_{k+1}R_k^*,
\qquad
R_k\coloneqq A\Sigma_{k+1}^{-1},
\]
and
\[
w_k^*\Pi_kw_k
=
\frac{2q}{\Delta_{k+1}}
\Bigl[
(\Delta_{k+1}+\eta v_{k+1})\,j_k
+
(\eta\tau_{k+1}+qy\,v_{k+1})\,|w_{k,2}|^2
\Bigr]
>0.
\]
The terminal term is also strictly positive:
\[
w_m^*P_mw_m>0
\]
in both the even and odd cases.
\end{proposition}

\begin{proof}[Proof sketch]
The recursion for \(S_k\) is
\[
S_k=I+R_kS_{k+1}R_k^*,
\]
so \(S_k\succ 0\). The formula for \(P_k\) is obtained by differentiation of the Riccati step,
and the displayed quadratic form is an exact computation. Its positivity follows from
\(q>0\), \(\Delta_{k+1}>0\), \(j_k>0\), and \(\tau_{k+1}>qy\). The terminal contribution is
computed explicitly; in the odd case one must differentiate the scalar tail factor
\[
e=x^2+y^2+b^2-\Lambda,
\]
and use that \(e>0\) at \(\Lambda=\lambda_*(x,y)\), since odd \(A_n\) has eigenvalue \(0\).
\end{proof}

\begin{theorem}[Strict monotonicity after subtracting the trivial \(y^2\)-term]
\label{thm:shifted-decreasing}
Let
\[
s_*(x,y)\coloneqq \lambda_*(x,y)-x^2-y^2.
\]
At every simple point of the lowest branch with \(y>0\), one has
\[
\partial_y s_*(x,y)<0.
\]
Equivalently,
\[
\partial_y\lambda_*(x,y)<2y.
\]
\end{theorem}

\begin{proof}
At a simple point of the lowest branch, choose the nullchain as above, with
\(\Lambda(y)=\lambda_*(x,y)\). Differentiating
\[
\Sigma_1(\Lambda(y),y)w_1(y)=0
\]
and pairing with \(w_1(y)\) gives
\[
\bigl(\Lambda'(y)-2y\bigr)\,w_1^*S_1w_1
=
-\,w_1^*P_1w_1.
\]
Now \(w_1^*S_1w_1>0\) and Proposition~\ref{prop:Ppositive} gives
\(w_1^*P_1w_1>0\). Hence \(\Lambda'(y)-2y<0\), i.e.
\[
\partial_y s_*(x,y)<0.
\]
\end{proof}

\begin{remark}
Theorem~\ref{thm:shifted-decreasing} is exact and general. However it is not by itself
sufficient for the full connectedness-equivalence result, because the desired statement is
\[
\partial_y\lambda_*(x,y)\ge 0
\qquad (y>0),
\]
which is stronger than \(\partial_y\lambda_*(x,y)<2y\).
\end{remark}

\section{What is proved for \(n=2,3\)}

\begin{theorem}\label{thm:n23}
For \(n=2\) and \(n=3\), one has
\[
|y|\longmapsto \lambda_*(x,y)
\]
nondecreasing for every fixed \(x\in\R\). Consequently, for \(n=2,3\),
\[
\sigma_\eps(A_n)\ \text{is connected}
\quad\Longleftrightarrow\quad
\sigma_\eps(A_n)\cap\R\ \text{is connected}
\qquad (\eps>0).
\]
\end{theorem}

\begin{proof}[Proof sketch]
For \(n=2\), \(\lambda_*(x,y)\) has a closed form and is manifestly nondecreasing in \(y^2\).
For \(n=3\), one reduces to a scalar equation in the variable \(\nu=x^2+y^2-\lambda_*(x,y)\)
and proves monotonicity in \(y^2\) by an explicit one-variable argument. Then
Lemma~\ref{lem:topological} applies.
\end{proof}

\section{The remaining gap}

The whole connectedness-equivalence problem has now been localized into one sign question.

\begin{conjecture}[Threshold-slope positivity]\label{conj:threshold}
For every \(n\ge 2\), every simple boundary point
\[
\Lambda=\lambda_*(x,y),
\qquad y>0,
\]
satisfies
\[
\partial_y \lambda_*(x,y)>0.
\]
Equivalently,
\[
w_1^*\,\partial_y\Sigma_1(\Lambda,y)\,w_1>0,
\]
where \(w_1\in\ker\Sigma_1(\Lambda,y)\setminus\{0\}\).
\end{conjecture}

\begin{remark}
Conjecture~\ref{conj:threshold} is the only missing theorem needed for the full
connectedness-equivalence result. Every exact reduction above is already proved.
\end{remark}

\begin{theorem}[Conditional completion of the real-slice criterion]
\label{thm:conditional}
Assume Conjecture~\ref{conj:threshold}. Then for every \(n\ge 2\) and every \(\eps>0\),
\[
\sigma_\eps(A_n)\ \text{is connected}
\quad\Longleftrightarrow\quad
\sigma_\eps(A_n)\cap\R\ \text{is connected}.
\]
\end{theorem}

\begin{proof}
Fix \(x\in\R\) and \(\eps>0\), and define
\[
F(y)\coloneqq \lambda_*(x,y)-\eps^2,
\qquad y\in\R.
\]
Then \(F\) is continuous and even, and
\[
F(y)\to +\infty
\qquad (|y|\to\infty),
\]
because \(M_n(x,y)\ge y^2I_n+O(|y|)\).

Now let \(y_0>0\) be a zero of \(F\). Then \(\eps^2=\lambda_*(x,y_0)\), so by
Conjecture~\ref{conj:threshold},
\[
F'(y_0)=\partial_y\lambda_*(x,y_0)>0.
\]
Hence every positive root is simple and is crossed from negative to positive as \(y\)
increases. Therefore \(F\) has at most one positive root. It follows that
\[
\{y\ge 0:\ F(y)<0\}
\]
is either empty or an interval of the form \([0,y_*)\). By evenness,
\[
\{y\in\R:\ F(y)<0\}
\]
is either empty or an interval containing \(0\). Thus every vertical slice of
\(\sigma_\eps(A_n)\) is either empty or an interval containing \(0\). The claim follows
from Lemma~\ref{lem:topological}.
\end{proof}

\section{Summary}

\begin{remark}
The exact progress to date is as follows.

\smallskip
\noindent
\textup{(1)} The full finite-\(n\) problem has been reduced to the \(2\times2\) parity Weyl
matrix \(\Sigma_1\).

\smallskip
\noindent
\textup{(2)} The tail Riccati dynamics are exact and identical in the even and odd cases,
up to the terminal seed.

\smallskip
\noindent
\textup{(3)} The block nullchain admits a projective upper half-plane invariant:
\[
r_k=\frac{w_{k,2}}{w_{k,1}}\in\mathbb H.
\]

\smallskip
\noindent
\textup{(4)} The associated cocycle \(P_k=2yS_k-\partial_y\Sigma_k\) is strictly positive
along the boundary nullchain, which proves
\[
\partial_y\bigl(\lambda_*(x,y)-x^2-y^2\bigr)<0.
\]

\smallskip
\noindent
\textup{(5)} The only unresolved step is the threshold-slope positivity
\[
\partial_y\lambda_*(x,y)>0
\qquad (y>0),
\]
which would finish the connectedness-equivalence theorem by
Theorem~\ref{thm:conditional}.
\end{remark}

\begin{proposition}[Takagi reduction and prefix identity]
Fix \(z=x+iy\) with \(y>0\), and let
\[
B(z):=zI-A_n,
\qquad
s:=\sigma_{\min}(B(z)).
\]
Choose a right singular vector \(u=(u_1,\dots,u_n)^T\neq 0\) in Takagi form:
\[
B(z)u=s\,J\overline{u},
\]
where \(J\) is the reversal matrix.

Define
\[
I_j:=\Im(\overline{u_j}u_{j+1}),
\qquad 1\le j\le n-1,
\qquad
I_0=I_n=0,
\]
and
\[
p_j:=-\Im(u_j u_{n+1-j}),
\qquad 1\le j\le n.
\]
Then the following exact identities hold:

\smallskip
\noindent
\textup{(i)} For each \(j=1,\dots,n\),
\[
s\,p_j
=
y|u_j|^2+a I_{j-1}-b I_j.
\]

\smallskip
\noindent
\textup{(ii)} For each \(L=1,\dots,n-1\), define
\[
\mathcal P_L
:=
y\sum_{j=1}^L |u_j|^2-(b-a)\sum_{j=1}^L I_j.
\]
Then
\[
\mathcal P_L
=
s\sum_{j=1}^L p_j + a I_L.
\]

\smallskip
\noindent
\textup{(iii)} If \(u_j\neq 0\) and \(u_{j-1}\neq 0\), and
\[
r_j:=\frac{u_{j+1}}{u_j},
\qquad
r_{j-1}:=\frac{u_j}{u_{j-1}},
\]
then
\[
\frac{u_j}{\overline{u_{n+1-j}}}
=
\frac{s}{x+iy-\dfrac{a}{r_{j-1}}-b r_j}.
\]
Hence
\[
p_j>0
\quad\Longleftrightarrow\quad
\Im\!\left(x+iy-\frac{a}{r_{j-1}}-b r_j\right)>0.
\]
\end{proposition}

\begin{proof}
The Takagi equation is
\[
(x+iy)u_j-a u_{j-1}-b u_{j+1}=s\,\overline{u_{n+1-j}},
\qquad
u_0=u_{n+1}=0.
\]
Multiplying by \(\overline{u_j}\) and taking imaginary parts gives
\[
s\,\bigl(-\Im(u_j u_{n+1-j})\bigr)
=
y|u_j|^2+a\,\Im(\overline{u_{j-1}}u_j)-b\,\Im(\overline{u_j}u_{j+1}),
\]
which is \textup{(i)}.

Summing \textup{(i)} from \(j=1\) to \(L\) yields
\[
s\sum_{j=1}^L p_j
=
y\sum_{j=1}^L|u_j|^2
+a\sum_{j=1}^L I_{j-1}
-b\sum_{j=1}^L I_j
=
y\sum_{j=1}^L|u_j|^2
-(b-a)\sum_{j=1}^{L-1}I_j
-b I_L.
\]
Rearranging gives \textup{(ii)}.

Finally, dividing the Takagi equation by \(\overline{u_{n+1-j}}\) gives
\[
s
=
\left(x+iy-\frac{a u_{j-1}}{u_j}-\frac{b u_{j+1}}{u_j}\right)\frac{u_j}{\overline{u_{n+1-j}}},
\]
which is \textup{(iii)}.
\end{proof}

\begin{conjecture}[Mirror-pair negativity]
With the notation above,
\[
p_j=-\Im(u_j u_{n+1-j})>0
\qquad\text{for every }j=1,\dots,n.
\]
Equivalently,
\[
\Im\!\left(x+iy-\frac{a}{r_{j-1}}-b r_j\right)>0
\]
for every \(j\).
\end{conjecture}

\begin{corollary}[Conditional completion of the connectedness-equivalence theorem]
Assume the mirror-pair negativity conjecture. Then for every \(n\ge 2\) and every \(\eps>0\),
\[
\sigma_\eps(A_n)\ \text{is connected}
\quad\Longleftrightarrow\quad
\sigma_\eps(A_n)\cap\R\ \text{is connected}.
\]
\end{corollary}

\begin{proof}
By the already established block theory, one has \(I_j>0\) for every \(j\), since these are
exactly the alternating within-block and cross-block currents.

By the proposition,
\[
\mathcal P_L
=
s\sum_{j=1}^L p_j + a I_L.
\]
Hence the conjecture implies \(\mathcal P_L>0\) for every \(L=1,\dots,n-1\). In particular,
for even \(L=2K\),
\[
\sum_{k=1}^K\Bigl(2y\|w_k\|^2-2(b-a)(j_k+h_k)\Bigr)=2\mathcal P_{2K}>0.
\]
Taking \(K=m-1\) gives positivity of the full interior threshold-derivative sum, and the
terminal threshold derivative was already proved positive in both the even and odd cases.
Therefore
\[
w_1^*\,\partial_y\Sigma_1\,w_1>0,
\]
so the boundary threshold \(\lambda_*(x,y)\) is strictly increasing in \(y>0\). Thus each
vertical slice of \(\sigma_\eps(A_n)\) is either empty or an interval containing \(0\), and the
connectedness equivalence follows from the vertical-slice topological lemma.
\end{proof}

\end{document}